\section{Related Work}
\label{sec:related}

\subsection{Federated Reinforcement Learning for Edge Orchestration}

Federated reinforcement learning (FRL) has emerged as a scalable and privacy-preserving solution for orchestrating services across edge networks~\cite{Wang2023_FedRLReview,Zhang2024_FedEdgeSurvey}. Unlike centralized RL, which requires global state observability, FRL enables distributed agents to learn coordinated policies while preserving local data privacy. Recent studies have adapted FRL for edge computing scenarios, introducing multi-agent coordination~\cite{Zhou2023_MARLFed}, policy distillation~\cite{Zhang2023_PolicyDistillFL}, and robustness under dynamic conditions~\cite{Nguyen2024_FedRLDynamic}. However, many of these frameworks rely on flat feature representations and do not encode the underlying network structure or the semantic context of service requests, limiting adaptability in heterogeneous and non-stationary edge environments.

\subsection{Semantic-Aware Scheduling and Service Placement}

Semantic representations provide high-level abstractions that enhance decision stability and generalization across orchestration tasks. Prior work has explored semantic embeddings for clustering, caching, and workload prediction in edge systems~\cite{Liu2024_SemanticEdge,Sun2024_SemanticScheduling}, as well as in personalized federated learning~\cite{Hu2023_PersonalizedFL} and offloading strategies~\cite{He2024_SemanticRL}. Nevertheless, reinforcement learning based orchestration methods have largely underutilized semantic similarity in service placement. While semantics have been employed for model personalization~\cite{Li2023_SemanticFL}, they remain underexplored in policy learning for distributed orchestration. In our evaluations, semantic clustering helps reduce service migrations and improves task continuity, contributing to system-level stability across diverse workload characteristics.

\subsection{GNN-Based Topology Encoding for Edge Intelligence}

Graph neural networks (GNNs) are increasingly used to model the structural relationships among edge nodes, enabling context-aware resource allocation and fault-tolerant decision-making~\cite{Kumar2023_GNNOrchestration,Chen2024_GNNTopologyEdge}. By encoding communication links and node states into graph representations, GCN-based models capture both local and global dependencies~\cite{Zhang2023_GNNResourceAlloc,Liang2024_TopoGNN}. GNNs have been applied to tasks such as multi-hop offloading and resilience-aware placement. However, most implementations assume centralized control or partial observability, and few operate under federated constraints. Furthermore, integration of GNNs into reinforcement learning pipelines is often limited to resource-centric optimization, without considering semantic consistency or communication efficiency. Our proposed approach addresses this by incorporating GCN-based topology modeling within a federated, PPO-driven orchestration loop.

\subsection{Hierarchical Federated Learning for Communication Efficiency}

Hierarchical federated learning (HFL) reduces synchronization overhead by grouping clients into subregions that perform local model updates before periodic global aggregation~\cite{Liu2023_HierFL,Zhao2023_HFLIoT,Tang2024_HFLV2X}. This model is especially suited for edge and IoT systems, where communication cost and energy efficiency are paramount. While HFL has shown promise in supervised learning, its use in reinforcement learning for real-time orchestration remains limited. Existing methods often neglect dynamic service behaviors or lack mechanisms for topology-informed coordination. In our work, we extend HFL principles to federated PPO and apply it in a two-tier control structure with local synchronization every 10 epochs and global synchronization every 50 epochs. This results in significant communication savings, reducing total communication volume to 0.65 MB over 1,000 orchestration steps compared to 105.0 MB for flat federated PPO, representing a 161 times reduction in communication overhead.


\subsection{Deep RL-Based Edge Offloading and Resource Allocation}

Recent work on deep reinforcement learning for edge computing has demonstrated significant advances in task offloading and resource allocation. ECO-SDIoT~\cite{zhu2023eco} employs Double DQN with software-defined networking for global-view optimization, achieving 20--60 ms task offloading latency in IoT environments. While effective for batch-oriented offloading, ECO-SDIoT lacks semantic awareness and operates without federated privacy guarantees. GFL-LFF~\cite{regan2024gfl} combines GNN with federated learning for privacy-preserving data aggregation in Industrial IoT, reporting 1.6 s latency for secure data sharing. Although GFL-LFF demonstrates the benefits of GNN integration with federated approaches, it focuses on data aggregation rather than real-time orchestration. FRPVC~\cite{qian2025frpvc} applies federated deep reinforcement learning with DDQN for cost-efficient video caching in energy-constrained mobile edge networks, optimizing cache hit rates through denoised autoencoders and MDP-based resource allocation.

While these methods demonstrate effective deep RL optimization for specific edge applications, they address narrow use cases (offloading, data aggregation, caching) without generalizing to semantic-aware task placement. Moreover, value-based methods (Double DQN, DDQN) can suffer from instability in high-dimensional action spaces compared to policy gradient approaches. FedSemGNN advances beyond these works by: (1) incorporating explicit semantic embeddings for application-aware orchestration, (2) employing hierarchical PPO for more stable policy learning than value-based methods, (3) using GNN topology encoding for efficient network representation, and (4) achieving sub-millisecond latency (0.36 ms) representing 55 to 4,444 times improvements over state-of-the-art methods while maintaining federated privacy guarantees.

\subsection{Positioning of Our Work}

To the best of our knowledge, no existing system jointly integrates semantic embeddings, GNN-based topology encoding, and hierarchical federated reinforcement learning within a unified orchestration framework. Our proposed \textbf{FedSemGNN} is the first to combine these dimensions in a federated setting, implemented on the EdgeSimPy simulator~\cite{edgesimpy}. The system supports detailed instrumentation across reward, latency, communication, power, and fidelity metrics. Evaluated over 1,000 orchestration steps on a heterogeneous edge infrastructure and extrapolated to scales of up to 10,000 nodes using validated mathematical models, FedSemGNN demonstrates substantial improvements compared to both self-implemented baselines (FlatFedPPO, HierFedPPO, HSQF, RandomPlacement) and state-of-the-art published methods (ECO-SDIoT, GFL-LFF, FRPVC). Key results include normalized reward of 0.91 (8.4 percent higher than FlatFedPPO), orchestration latency of 0.36 ms (315 times faster than FlatFedPPO, 55 to 166 times faster than ECO-SDIoT, and 4,444 times faster than GFL-LFF), 100 percent semantic fidelity with zero migrations, communication efficiency of 1.394 reward per MB (76 times better than ECO-SDIoT), and energy consumption of 72.1 W (36.2 times lower than FlatFedPPO). These findings highlight the architectural potential of semantic and topological modeling in federated edge orchestration.
